%% Example CV - 王鹏生
%% Comprehensive overview of all competencies, experience, and achievements.
%% Use as a master reference when tailoring role-specific CVs.
%%
%% CV language: 中文（/setup 确认）。章节标题与样板行均为中文。
%% Compile with: lualatex main_example.tex
%% (pdflatex often fails on modern MiKTeX with fontawesome5 font-expansion errors.)

\documentclass[11pt,a4paper,sans]{moderncv}
\moderncvstyle{banking}
\moderncvcolor{blue}

% Force the name and section headings to render in moderncv blue (color1).
% Default banking leaves them black: moderncvstylebanking.sty's \colorlet
% copies (not aliases) the pre-scheme accent colour, so the name colours are
% frozen before \moderncvcolor runs. Re-let them after. \namefont is the hook
% every name-style macro routes through, so this also works on moderncv 2.3.1
% (Debian/Ubuntu apt), which has no \firstnamestyle/\lastnamestyle at all.
\renewcommand*{\namefont}{\fontsize{34}{36}\bfseries\upshape}
\colorlet{firstnamecolor}{color1}
\colorlet{lastnamecolor}{color1}
\colorlet{namecolor}{color1}
\renewcommand*{\sectionstyle}[1]{{\sectionfont\color{color1}#1}}

% 中文支持：moderncv 默认只有拉丁字体，缺 CJK 设置时中文会缺字或显示为方框。
% 字体名需与系统已装字体一致（Windows: Microsoft YaHei / SimSun / SimHei；
% 跨平台: Noto Sans CJK SC）。编译仍用 lualatex。
\usepackage{ctex}
\setCJKmainfont{SimSun}
\setCJKsansfont{Microsoft YaHei}

\usepackage[utf8]{inputenc}
% moderncv loads hyperref itself in an \AtEndPreamble hook, so \hypersetup
% must go in an \AtEndPreamble of our own: on moderncv < 2.4 a top-level
% \usepackage{hyperref} clashes with the class's own
% \RequirePackage[unicode]{hyperref}. From 2.4.0 the class passes its options
% through \PassOptionsToPackage instead, which is what removes that clash.
\AtEndPreamble{\hypersetup{
    colorlinks=true,
    linkcolor=blue,
    filecolor=magenta,
    urlcolor=blue,
    pdftitle={王鹏生 - CV},
    % Keep pdfpagemode=UseNone: this block runs after moderncv's own
    % \AtEndPreamble (moderncv.cls sets pdfpagemode there), so a FullScreen
    % value here would win and open every CV in fullscreen presentation mode.
    pdfpagemode=UseNone,
}}
\usepackage[scale=0.80]{geometry}
\usepackage{import}
\usepackage{needspace}

% personal data
\name{王}{鹏生}
% If you have no address to list, DELETE this whole line. \address{}{}{} fails
% with "There's no line here to end" on every moderncv version.
\address{北京}{}{}
\phone[mobile]{18721695636}
\email{wpsseu@163.com}
% 未提供 LinkedIn / GitHub，故不输出 \extrainfo 行（占位 URL 会进编译产物）。

\begin{document}

\makecvtitle

% ============================================================
%     PROFILE STATEMENT
% ============================================================

\vspace{6pt}
\small{具备人工智能与汽车复合能力的测试与质量负责人，13 年+汽车行业经验、6 年测试开发管理经验，曾任职头部 AI 企业与头部车企。深耕智能座舱传统 AI 功能与大模型评测方向，主导交付国内外 10 余家车企的 30 余个智舱量产及 POC 项目，参与车载 AI 行业标准制定。熟悉智能座舱从软件供应到整车落地的全链路测试流程，可独立从 0 到 1 搭建自动化测试框架以及传统模型、大模型评测体系。获 2024 年商汤集团年度杰出个人奖。}

% ============================================================
%     CORE COMPETENCIES
% ============================================================

\section{核心能力}
\vspace{1pt}
\begin{itemize}

\item \textbf{AI 与大模型评测}: 精通多类 AI 评测方法，熟练运用 PyTorch 开展模型性能优化与压力测试；可搭建 VFM、VLM、LLM、Agent、多模态模型测试数据集，实现 Prompt 批量测试、自动评分、Badcase 分析及回归验证；围绕准确率、召回率、幻觉率、时延、鲁棒性构建评测指标体系，支撑 SFT 微调版本对比。

\item \textbf{智能座舱全链路测试}: 覆盖需求评审、用例设计、集成、功能、性能精度、稳定性、可靠性、安全测试以及台架与实车测试；可完成模型、SDK、应用、硬件及 HMI 显示各维度测试，保障量产交付。

\item \textbf{测试开发与自动化}: 熟练使用 Python、C++、Java、TypeScript、Shell，可独立设计开发自动化测试框架、AI 测试平台与数据分析工具；掌握 Selenium、Postman、JMeter，具备接口自动化与 CI/CD 落地经验。

\item \textbf{测试管理与量产交付}: 承担量产项目测试管理与策略制定，服务国内外十余家 OEM 及 Tier One 客户（含德国、日本及国内主流车企），落地近 30 款车型量产交付。

\item \textbf{标准与法规}: 深度解读车载质量标准，推动产品通过信息安全、A-SPICE L2/L3、MLE 认证；参与车用 AI 专项、自动驾驶接管能力监测、汽车座舱功能评价技术规范编制。

\end{itemize}

% ============================================================
%     PROFESSIONAL EXPERIENCE
% ============================================================

\section{工作经历}
\vspace{3pt}
\begin{itemize}

% --- Most Recent Role ---
\item{\cventry{2026.07-至今}{测试开发}{一汽软件}{}{}{\vspace{1pt}
\begin{itemize}
    \item {牵头智能座舱整体测试体系建设，全权负责屏幕黑白屏、花屏、闪屏等 HMI 显示类专项测试方案设计与落地}
    \item {搭建包含用例设计、缺陷识别、根因定位的端到端测试能力，协同研发定位显示异常并推动缺陷闭环验证}
    \item {负责自动化测试框架搭建：用例批量执行、测试结果自动采集、异常自动识别}
    \item {输出标准化测试台架方案并落地远程云真机平台，扩大自动化回归覆盖，提升跨地域测试执行效率}
\end{itemize}}}

\vspace{3pt}

% --- Previous Role ---
\item{\cventry{2022.12-2026.06}{测试经理}{上海商汤智能科技有限公司}{}{}{\vspace{1pt}
\begin{itemize}
    \item {团队管理与量产交付：覆盖 DMS、OMS、Face ID、哨兵模式等功能，服务国内外十余家 OEM 及 Tier One 客户，落地近 30 款车型量产项目}
    \item {大模型产品预研与测试体系搭建：牵头 LLM、VLM 模型侧至 Agent、APP 应用侧的全链路多维度测试策略规划，带领团队构建百万级测试数据集}
    \item {支撑 POI 推荐、车载说明书、AI 壁纸生成、健康助手、遗留物监测与儿童关怀 Agent 等大模型产品落地，累计反馈 400+ 条 Badcase，获客户感谢信}
    \item {推动产品通过信息安全、A-SPICE L2/L3、MLE 认证；多次代表公司出席汽车行业标准会议并参与 3 项技术规范编制}
\end{itemize}}}

\vspace{3pt}

% --- Earlier Role ---
\item{\cventry{2020.12-2022.11}{测试开发}{华为苏州研究所}{}{}{\vspace{1pt}
\begin{itemize}
    \item {主导搭建企业级自动化测试框架，基于 Python、Shell 编写 1000+ 测试脚本，推动团队自动化率从 60\% 提升至 90\%}
    \item {实现故障注入、版本升级、性能压测等高难度专项测试全流程自动化，建立自动化用例库与持续回归机制}
    \item {独立研发自动化测试工具与环境自动部署工具，测试环境部署耗时从 2 小时压缩至 30 分钟}
\end{itemize}}}

\vspace{3pt}

\item{\cventry{2013.07-2020.11}{机械结构工程师}{上海翼锐汽车科技有限公司}{}{}{\vspace{1pt}
\begin{itemize}
    \item {负责底盘零件结构设计：主导开发多款车型车轮、随车工具、紧固件等零件，7 年整车研发经验奠定汽车行业流程与质量体系基础}
\end{itemize}}}

\end{itemize}

% ============================================================
%     EDUCATION
% ============================================================

\section{教育背景}
\vspace{1pt}
\begin{itemize}

\item{\cventry{2018-2021}{硕士（非全日制），计算机科学与技术}{上海交通大学}{}{}{\vspace{1pt}
与华为苏州研究所任职期重叠，属非全日制在读。
}}

\vspace{3pt}

\item{\cventry{2009-2013}{本科，机械制造与自动化}{东南大学}{}{}{\vspace{1pt}
机械制造与自动化。
}}

\end{itemize}

% ============================================================
%     LANGUAGES
% ============================================================

\section{语言}
\vspace{1pt}
\begin{itemize}
\item {中文（母语），英语（CET-6，可进行技术会议交流）。}
\end{itemize}

% ============================================================
%     HONORS AND AWARDS
% ============================================================

\section{荣誉奖项}
\vspace{1pt}
\begin{itemize}
\item{\textbf{2024 年商汤集团年度杰出个人奖} -- 商汤集团 (2024)。}
\item{\textbf{客户感谢信} -- 大模型项目交付获客户高度赞扬。}
\item{\textbf{参与编制 3 项车载 AI / 座舱技术规范} -- 车用 AI 专项、自动驾驶接管能力监测、汽车座舱功能评价。}
\end{itemize}

% ============================================================
%     REFERENCES
% ============================================================

\section{推荐人}
\vspace{1pt}
\begin{itemize}
\item{如有需要可提供推荐人信息。}
\end{itemize}

\end{document}
